\documentclass[12pt]{article}
\usepackage{algorithmic}
\usepackage{algorithm}
\usepackage{alltt}
\usepackage{amsbsy}
\usepackage{amscd}
\usepackage{amsfonts}
\usepackage{amsmath}
\usepackage{amssymb}
\usepackage{amsthm}
\usepackage{array}
\usepackage{booktabs}
\usepackage{epsfig}
\usepackage{euscript}
\usepackage{fancyhdr}
\usepackage{float}
\usepackage{fullpage}
\usepackage{graphicx}
\usepackage{ifthen}
\usepackage{multirow}
\usepackage{supertabular}
\usepackage{times}
\usepackage{verbatim}
\usepackage{xfrac}
\usepackage{xspace}
\usepackage[colorlinks]{hyperref}
\usepackage[numbers,sort&compress]{natbib}
\usepackage{subfig}
\usepackage[normalem]{ulem}
\usepackage{url}

\usepackage{boldtensors}

\usepackage[T1]{fontenc}
\usepackage{lmodern}

\newcommand{\F}{~F}
\newcommand{\C}{~C}
\renewcommand{\P}{~P}
\newcommand{\eps}{~\epsilon}
\newcommand{\epsbar}{\bar{~\epsilon}}
\newcommand{\edev}{~e}
\newcommand{\ebar}{\bar{~e}}
\newcommand{\btau}{~\tau}
\newcommand{\btaubar}{\bar{~\tau}}
\newcommand{\sdev}{~s}
\newcommand{\sbar}{\bar{~s}}
\newcommand{\thbar}{\bar{\theta}}
\newcommand{\pbar}{\bar{p}}
\newcommand{\Reals}{\mathbb{R}}
\newcommand{\Syms}{\mathbb{S}}
\newcommand{\supvol}{^{\mathrm{vol}}}
\newcommand{\supdev}{^{\mathrm{dev}}}
\newcommand{\Symsd}{\mathbb{S}\supdev}
\newcommand{\norm}[1]{\left\lVert {#1} \right\rVert}
\newcommand{\abs}[1]{\left\lvert {#1} \right\rvert}
\DeclareMathOperator{\trace}{tr}
\DeclareMathOperator{\dev}{dev}
\DeclareMathOperator{\vol}{vol}
\DeclareMathOperator{\Div}{Div}
\DeclareMathOperator{\Grad}{Grad}
\DeclareMathOperator{\Curl}{Curl}
\DeclareMathOperator{\sym}{sym}
\newcommand*\diff{\mathop{}\!\mathrm{d}}
\newcommand{\transpose}[1]{{#1}^{\mathrm{T}}}

\begin{document}

\setlength{\headheight}{15pt}
\headsep = 4pt
\pagestyle{fancyplain}

\lhead
[\fancyplain{}{\emph{A. Mota \& J. Foulk}}]
{\fancyplain{}{\emph{A. Mota \& J. Foulk}}}

\rhead
[\fancyplain{}{\emph{SEC5L Tetrahedron}}]
{\fancyplain{}{\emph{SEC5L Tetrahedron}}}

\title{
  SEC5L: A Five-Field Mixed Tetrahedral Finite Element on the
  Symmetric Elasticity Complex with Logarithmic Strain}

\author{
  Alejandro Mota$^1$\thanks{Email: amota@sandia.gov},
  James W. Foulk {III}$^1$\thanks{Email: jwfoulk@sandia.gov}\\
  \\
  $^1$Mechanics of Materials Department\\
  Sandia National Laboratories\\
  Livermore CA 94550, USA\\
}

\date{\today}

\maketitle

\begin{abstract}

  We propose the SEC5L element, a tetrahedral finite element built
  on the symmetric elasticity complex (SEC) with 5 independent
  discrete fields and a logarithmic strain measure. The five fields
  are displacement, volumetric log-strain, deviatoric log-strain,
  hydrostatic pressure, and deviatoric Kirchhoff stress. The
  formulation is structured so that all
  finite-deformation nonlinearity is absorbed into the kinematic
  relation~$\eps = \tfrac{1}{2}\log(\transpose{\F}\F)$, while the
  constitutive relation is algebraically linear in log-strain
  space. The orthogonal volumetric--deviatoric decomposition
  $\eps = \tfrac{1}{3}\theta~I + \edev$, with $\theta := \trace\eps$
  and $\edev := \dev\eps$, and $\btau = p~I + \sdev$, with
  $p := \tfrac{1}{3}\trace\btau$ and $\sdev := \dev\btau$, is native
  to the Hencky measure, which
  decouples the $J_2$ plastic return map (pointwise on $\edev$)
  from the volumetric response and exposes a natural handle for stiff
  bulk-wave regimes. The independent strain and stress fields remove
  the cross-space constitutive evaluation that compromises
  stabilized displacement formulations, and all four auxiliary fields
  admit exact $L^2$ projections from $~\varphi$ with mass matrices
  that are constant in the reference configuration and
  pre-factorizable, reducing the global nonlinear solve to a
  displacement-sized problem. The formulation is intended as an
  alternative to the composite tetrahedron of \citet{Ostien.etal:2016}
  and its five-field reformulation \citep{Foulk.etal:2021}, targeting
  near-incompressible finite-deformation elasticity, isochoric
  plasticity, and weak-shock regimes with large bulk-to-shear ratio.
  \\
  \\
  Keywords: SEC5L, Hu-Washizu, tetrahedron, mixed formulation,
  logarithmic strain, Kirchhoff stress, symmetric elasticity complex,
  plasticity.

\end{abstract}

\section{Introduction}
\label{sec:introduction}

Tetrahedral finite elements are the natural choice for discretizing
solid bodies with complex geometry because they admit robust
automatic meshing. Their performance in problems involving
near-incompressible materials, isochoric plasticity, and stiff
bulk-wave regimes has, however, long been the limiting factor in
their adoption for large-deformation Lagrangian solid mechanics
\citep{Thoutireddy.etal:2002, Puso.Solberg:2006, Irving.etal:2007,
Danielson:2014, Scovazzi:2015}. Stabilized displacement methods
\citep{Klaas:1999, Maniatty:2002, Scovazzi:2015} and assumed-gradient
schemes such as the composite tetrahedron of \citet{Ostien.etal:2016}
have made significant progress. The composite tetrahedron with
mean-dilatation patch targets isochoric plasticity and
near-incompressible response through a three-field Hu-Washizu
functional with element-local discontinuous strain and stress and
static condensation. The subsequent five-field reformulation of
\citet{Foulk.etal:2021} adds independent Jacobian and pressure fields
to address a deleterious soft mode that arises in constant-pressure
formulations.

The present formulation departs from the composite-tetrahedron
lineage in three specific ways. First, the independent strain and
stress fields are discretized in \emph{conforming} symmetric-tensor
spaces compatible with the symmetric elasticity complex
\citep{Arnold.Falk.Winther:2006a, Arnold.Falk.Winther:2010}, giving
pointwise Cauchy symmetry of the discrete stress and normal-traction
continuity across element faces as structural properties of the
finite element spaces rather than of the constitutive law. Second,
the finite-deformation measure is the Hencky log-strain
$\eps := \tfrac{1}{2}\log(\transpose{\F}\F)$, with its work-conjugate
Kirchhoff stress $\btau$; the Hencky decomposition
$\eps = \tfrac{1}{3}(\log J)~I + \edev$ is orthogonal in the
Frobenius inner product and separates naturally into independent
volumetric and deviatoric discrete fields. Third, the constitutive
law is \emph{algebraically linear} in $(\eps, \btau)$ for isotropic
hyperelastic response, so the radial return map for $J_2$ plasticity
reduces exactly to the small-strain algorithm of \citet{Simo:1992}
acting pointwise on the deviatoric strain at quadrature.

The combination of these choices yields a formulation in which the
four auxiliary fields are obtained from the displacement field by
$L^2$ projections with mass matrices that are constant in the
reference configuration. Pre-factorization at setup reduces each
Newton iteration to a displacement-sized nonlinear solve, while the
auxiliary fields preserve the spatial-discretization correctness
properties (symmetry, traction continuity) through their conforming
finite element spaces. Dynamic and quasi-static regimes use the same
spatial formulation; the mass matrix enters only through the
displacement block.


\section{Continuum Formulation}
\label{sec:continuum}

Consider a body $B \subset \Reals^3$ undergoing a motion described
by $~x = ~\varphi(~X, t): B \times [t_1, t_2] \rightarrow \Reals^3$,
with deformation gradient $\F := \Grad ~\varphi$ and Jacobian
$J := \det \F$. The boundary $\partial B$ with outward unit normal
$~N$ is partitioned into a displacement boundary
$\partial_\varphi B$ and a traction boundary $\partial_T B$.
Prescribed boundary data are
$~\chi : \partial_\varphi B \times [t_1, t_2] \rightarrow \Reals^3$
and $~T : \partial_T B \times [t_1, t_2] \rightarrow \Reals^3$.
Let $R : B \rightarrow \Reals$ be the reference mass density and
$~B : B \times [t_1, t_2] \rightarrow \Reals^3$ the body force per
unit mass.

\subsection{Hencky Log-Strain and Kirchhoff Stress}
\label{sec:hencky}

Let $\C := \transpose{\F}\F$ be the right Cauchy--Green tensor. The
material Hencky log-strain is
\begin{equation}
  \label{eq:hencky}
  \eps := \tfrac{1}{2} \log \C \in \Syms,
\end{equation}
evaluated numerically by an inverse scaling-and-squaring Pad\'e
approximant applied to $\C$. Its orthogonal
volumetric--deviatoric decomposition is
\begin{equation}
  \label{eq:voldev-strain}
  \eps = \tfrac{1}{3}\,\theta~I + \edev,
  \qquad
  \theta := \trace \eps = \log J,
  \qquad
  \edev := \dev \eps.
\end{equation}
The work-conjugate stress measure is the Kirchhoff stress
$\btau$, symmetric and defined such that $\btau : \dot{\eps}$ equals the
stress power per unit reference volume for isotropic material
response. Its decomposition mirrors~\eqref{eq:voldev-strain}:
\begin{equation}
  \label{eq:voldev-stress}
  \btau = p~I + \sdev,
  \qquad
  p := \tfrac{1}{3} \trace \btau,
  \qquad
  \sdev := \dev \btau.
\end{equation}

\subsection{Stored Energy and Constitutive Law}
\label{sec:stored-energy}

For isotropic hyperelasticity, the Helmholtz free-energy density is
additively decomposable in the Hencky measure:
\begin{equation}
  \label{eq:W-split}
  A(\eps) = W\supvol(\theta) + W\supdev(\edev),
  \qquad
  W\supvol(\theta) = \frac{\kappa}{2} \theta^2,
  \qquad
  W\supdev(\edev) = \mu \, \norm{\edev}^2,
\end{equation}
with $\mu$ the shear modulus and $\kappa$ the bulk modulus. The
constitutive law is \emph{algebraically linear}:
\begin{equation}
  \label{eq:constitutive}
  p = \frac{\partial W\supvol}{\partial \theta} = \kappa \theta,
  \qquad
  \sdev = \frac{\partial W\supdev}{\partial \edev} = 2\mu \edev.
\end{equation}
Extensions to elasto-plasticity replace~\eqref{eq:constitutive}$_2$ by
\begin{equation}
  \label{eq:constitutive-plastic}
  \sdev = 2\mu \left( \edev - \edev^{\mathrm{p}} \right),
\end{equation}
with $\edev^{\mathrm{p}}$ the deviatoric plastic log-strain
evolved by a $J_2$ radial return map described
in~\S\ref{sec:plasticity}.

\subsection{Five-Field Hu-Washizu Functional}
\label{sec:hw-functional}

Introduce independent strain fields $\thbar$ and $\ebar$, and
enforce their equality with the pointwise kinematic fields
$\theta(~\varphi)$ and $\edev(~\varphi)$ weakly through Lagrange
multipliers $\pbar$ and $\sbar$. At a stationary point the
multipliers coincide with the work-conjugate stress fields: the
hydrostatic pressure and the deviatoric Kirchhoff stress,
respectively. The Hu-Washizu functional is
\begin{equation}
  \label{eq:hw}
  \begin{split}
    \Phi[~\varphi, \thbar, \ebar, \pbar, \sbar] & :=
    \int_B W\supvol(\thbar) \diff V
    +
    \int_B W\supdev(\ebar) \diff V
    \\
    & \quad
    +
    \int_B \pbar \left( \theta - \thbar \right) \diff V
    +
    \int_B \sbar : \left( \edev - \ebar \right) \diff V
    \\
    & \quad
    -
    \int_B R ~B \cdot ~\varphi \diff V
    -
    \int_{\partial_T B} ~T \cdot ~\varphi \diff S,
  \end{split}
\end{equation}
with $\theta$ and $\edev$ understood as pointwise functions
of $~\varphi$ through~\eqref{eq:hencky}--\eqref{eq:voldev-strain}.

\subsection{Variations and Euler-Lagrange Equations}
\label{sec:variations}

Assume $~\varphi \in U := (W^1_2(B))^3$,
$\thbar, \pbar \in Q := L^2(B)$,
$\ebar \in V_e := H(\Curl; B; \Symsd)$, and
$\sbar \in V_s := H(\Div; B; \Symsd)$, where $\Symsd$ denotes the space
of symmetric trace-free second-order tensors. Test functions are
$~\psi \in U$ with $~\psi = \mathbf{0}$ on
$\partial_\varphi B$, $\vartheta, \pi \in Q$, $~\eta \in V_e$, and
$~\xi \in V_s$. The variations of~\eqref{eq:hw} are
\begin{equation}
  \label{eq:variations}
  \begin{split}
    D_{~\varphi}\Phi(~\psi)
    & =
    \int_B \pbar \, \frac{\partial \theta}{\partial ~\varphi}[~\psi] \diff V
    +
    \int_B \sbar : \frac{\partial \edev}{\partial ~\varphi}[~\psi] \diff V
    \\
    & \quad
    -
    \int_B R ~B \cdot ~\psi \diff V
    -
    \int_{\partial_T B} ~T \cdot ~\psi \diff S
    = 0,
    \\
    D_{\thbar}\Phi(\vartheta)
    & =
    \int_B \left( \frac{\partial W\supvol}{\partial \thbar} - \pbar \right) \vartheta \diff V
    = 0,
    \\
    D_{\ebar}\Phi(~\eta)
    & =
    \int_B \left( \frac{\partial W\supdev}{\partial \ebar} - \sbar \right) : ~\eta \diff V
    = 0,
    \\
    D_{\pbar}\Phi(\pi)
    & =
    \int_B \left( \theta - \thbar \right) \pi \diff V
    = 0,
    \\
    D_{\sbar}\Phi(~\xi)
    & =
    \int_B \left( \edev - \ebar \right) : ~\xi \diff V
    = 0.
  \end{split}
\end{equation}
The two kinematic directional derivatives appearing
in~\eqref{eq:variations}$_1$ admit the explicit forms
\begin{equation}
  \label{eq:kinematic-derivatives}
  \frac{\partial \theta}{\partial ~\varphi}[~\psi]
  =
  \F^{-\mathrm{T}} : \Grad ~\psi,
  \qquad
  \frac{\partial \edev}{\partial ~\varphi}[~\psi]
  =
  \mathbb{L}\supdev(\C) : \sym\!\left(\transpose{\F}\Grad ~\psi\right),
\end{equation}
where $\mathbb{L}(\C) := \partial \log \C / \partial \C$ is the
fourth-order Fr\'echet derivative of the matrix logarithm and
$\mathbb{L}\supdev(\C) := \mathbb{L}(\C) - \tfrac{1}{3}\, ~I \otimes \C^{-1}$
its deviatoric part. The volumetric expression is algorithm-free
thanks to Jacobi's formula, while the deviatoric one requires a
dedicated evaluation of $\mathbb{L}(\C)$; the derivation and a
quadrature-point evaluation strategy are given
in~\ref{app:variations}.

The corresponding Euler-Lagrange equations are
\begin{equation}
  \label{eq:euler-lagrange}
  \begin{split}
    \Div \P + R ~B & = \mathbf{0} \quad \text{in} \quad B,
    \quad
    \P ~N = ~T \quad \text{on} \quad \partial_T B,
    \\
    \pbar & = \partial W\supvol / \partial \thbar = \kappa \thbar \quad \text{in} \quad B,
    \\
    \sbar & = \partial W\supdev / \partial \ebar = 2\mu \ebar \quad \text{in} \quad B,
    \\
    \thbar & = \log J \quad \text{in} \quad B,
    \\
    \ebar & = \dev \left[ \tfrac{1}{2} \log (\transpose{\F}\F) \right] \quad \text{in} \quad B,
  \end{split}
\end{equation}
where $\P$ is the first Piola-Kirchhoff stress reconstructed from
the independent fields as
\begin{equation}
  \label{eq:P-reconstruction}
  \P = (\pbar ~I + \sbar) \, \F^{-\mathrm{T}}
  \quad\text{(equivalently } \P = \btau \, \F^{-\mathrm{T}} \text{ with } \btau = \pbar ~I + \sbar \text{)},
\end{equation}
with $\btau = J ~\sigma$ the spatial Kirchhoff stress. Equivalently, in
a formulation organized around the Kirchhoff measure, $\P$
is obtained by the push-forward appropriate to the chosen
work-conjugate pair.


\section{Discretization}
\label{sec:discretization}

\subsection{Finite Element Spaces}
\label{sec:fe-spaces}

Introduce discretizations on a simplicial mesh of $B$:
\begin{equation}
  \label{eq:basis-functions}
  \begin{split}
    ~\varphi(~X) & := N_a(~X) ~\varphi_a \in U_h \subset U,
    \\
    \thbar(~X) & := \mu_m(~X) \, \thbar_m \in Q_{\theta,h} \subset Q,
    \\
    \ebar(~X) & := ~\lambda^e_\alpha(~X) \, \ebar_\alpha \in V_{e,h} \subset V_e,
    \\
    \pbar(~X) & := \mu_n(~X) \, \pbar_n \in Q_{p,h} \subset Q,
    \\
    \sbar(~X) & := ~\lambda^s_\beta(~X) \, \sbar_\beta \in V_{s,h} \subset V_s.
  \end{split}
\end{equation}
At polynomial degree $k \geq 2$, a stable choice on tetrahedra is:
\begin{itemize}
  \item $N_a$: continuous barycentric polynomial of degree $k$ on the tetrahedron for the displacement;
  \item $\mu_m, \mu_n$: discontinuous $L^2$ scalar of degree $k-1$
    for $\thbar$ and $\pbar$;
  \item $~\lambda^e_\alpha$: symmetric-tensor-valued basis for
    $H(\Curl; \Symsd)$ of degree $k-1$ (e.g.\ Nédélec-type rows with
    symmetry and trace-free constraints enforced algebraically);
  \item $~\lambda^s_\beta$: symmetric-tensor-valued basis for
    $H(\Div; \Symsd)$ of degree $k-1$, using the weakly-symmetric
    construction of \citet{Arnold.Falk.Winther:2007} with a skew
    multiplier in $L^2(B; \mathbb{K})$, or the strongly-symmetric
    construction of \citet{Arnold.Awanou.Winther:2008}.
\end{itemize}
Under these choices the discrete formulation satisfies the relevant
inf-sup stability conditions for the coupled Hu-Washizu system, with
constants independent of the bulk-to-shear ratio up to the range
$\kappa/\mu \lesssim 10^4$ relevant to the target problem class. Full
incompressibility ($\kappa \rightarrow \infty$) is outside the scope
of the present formulation.

Two aspects of the space choice deserve comment: why $\thbar$ and
$\pbar$ are \emph{allowed} to be discontinuous while $\ebar$ and
$\sbar$ are not, and why the polynomial degrees in the four
auxiliary slots are forced to $k-1$.

\textit{Integrability versus stability.} The Sobolev regularity
required of each auxiliary field must be addressed in two stages.
First, the HW functional~\eqref{eq:hw} and its
variations~\eqref{eq:variations} must be well-defined: none of the
four auxiliary fields appears under a spatial derivative, so
$L^2(B)$ is sufficient for integrability of every term (the volumetric
energy $W\supvol(\thbar)$ depends on values only; the Lagrange
pairings $\int \pbar(\theta - \thbar)$ and $\int \sbar:(\edev -
\ebar)$ carry no derivatives on the auxiliary fields; and in the
momentum variation~\eqref{eq:variations}$_1$ the gradient acts on
the test function $~\psi$, not on $\pbar$ or $\sbar$). Integrability
alone would therefore admit $L^2$ for all four. Second, however, the
coupled saddle-point system must be \emph{inf-sup stable} uniformly
in $h$ and in $\kappa/\mu$, and stability is a stronger requirement
than integrability.

Stability is inherited from the symmetric elasticity complex
\begin{equation}
  \label{eq:elasticity-complex}
  H^1(B; \mathbb{R}^3)
  \;\xrightarrow{\mathrm{sym}\,\Grad}\;
  H(\Curl; B; \Syms)
  \;\xrightarrow{\mathrm{inc}}\;
  H(\Div; B; \Syms)
  \;\xrightarrow{\Div}\;
  L^2(B; \mathbb{R}^3)
\end{equation}
\citep{Arnold.Falk.Winther:2006a,Arnold.Falk.Winther:2010}, whose
discrete counterpart admits cochain projections and a commuting
diagram when, and only when, each field sits in the conformity
dictated by its slot in the sequence. Consider, for reference, a
three-field HW formulation in the primal fields
$(~\varphi, \epsbar, \btaubar)$, where $\epsbar \in H(\Curl; \Syms)$
and $\btaubar \in H(\Div; \Syms)$ are independent symmetric-tensor
fields constrained only weakly to match $\eps(~\varphi)$ and
$\btau(\epsbar)$ through Lagrange-multiplier couplings. Even though
$\epsbar$ and $\btaubar$ are themselves independent, their
volumetric traces $\theta = \tfrac{1}{3}\trace \epsbar$ and $p =
\tfrac{1}{3}\trace \btaubar$ are \emph{embedded scalar components}
of globally-coupled tensor fields: the tangential-continuity
requirement on $\epsbar$ and the normal-continuity requirement on
$\btaubar$ prevent these scalar components from being chosen
independently in discontinuous $L^2$.

In the five-field decomposition, by contrast, $\thbar$, $\ebar$,
$\pbar$, $\sbar$ are introduced as \emph{independent} auxiliary
fields, related to the trace and trace-free parts of $\eps$ and
$\btau$ only weakly through the Lagrange-multiplier constraints.
Peeling the trace parts out of the tensor fields severs their
tensorial coupling: each auxiliary field is free to sit in
whatever space is compatible with its own weak form, subject only
to the inf-sup requirement of the coupled system. The stability of the coupled system then factors along the
orthogonal volumetric--deviatoric split. The deviatoric part
$(~\varphi, \ebar, \sbar)$ is governed by a symmetric elasticity
complex restricted to the trace-free subspace $\Symsd$, which
dictates $\ebar \in H(\Curl; \Symsd)$ and $\sbar \in H(\Div;
\Symsd)$; here $L^2$ conformity is insufficient and admits
spurious stress modes (the symmetric-tensor analog of
$P_1$-$P_0$ checkerboarding). The volumetric part $(~\varphi,
\thbar, \pbar)$ is governed by a Stokes-type mixed sub-problem
(detailed below), whose stability is classical and is attained
with discontinuous $L^2$ scalars of degree $k-1$. Discontinuous
$L^2$ bases are therefore admissible for $\thbar$ and $\pbar$, and
inadmissible for $\ebar$ and $\sbar$: not because the scalars stand
outside the complex algebraically, but because treating them as
independent fields in a Stokes-like sub-problem is what permits,
and is then stabilized by, the weaker $L^2$ conformity.

\textit{The volumetric sub-problem as a Stokes-type system.} By
``Stokes-type mixed sub-problem'' we mean a saddle-point problem
with a vector primal field in $H^1$, a scalar multiplier in $L^2$,
and an inf-sup pairing of the form $\int q \,\mathcal{D} v$ where
$\mathcal{D}$ is a first-order scalar differential operator on the
primal. The classical Stokes problem
\begin{equation}
  \label{eq:stokes-classical}
  -\mu \Delta u + \Grad p = f,
  \qquad
  \Div u = 0,
\end{equation}
is the prototype: $u \in H^1(B; \mathbb{R}^3)$, $p \in L^2(B)$, and
$p$ enforces $\Div u = 0$ through the pairing $\int p \, \Div v$.
Its stability is the Babu\v{s}ka--Brezzi condition
\begin{equation}
  \label{eq:babuska-brezzi}
  \inf_{q \in L^2}\;\sup_{v \in H^1}\;
  \frac{\int_B q \, \Div v \diff V}
       {\norm{v}_{H^1} \norm{q}_{L^2}}
  \;\geq\; \beta > 0,
\end{equation}
and the discrete pairs satisfying its analog uniformly in $h$
(Taylor--Hood $P_k$--$P_{k-1}$ continuous, $P_k$--$P_{k-1}$
discontinuous, MINI, Crouzeix--Raviart) are classical.

The volumetric sub-problem $(~\varphi, \thbar, \pbar)$ has
the same skeleton, with three substitutions summarized in
Table~\ref{tab:stokes-correspondence}.
\begin{table}[h]
  \centering
  \begin{tabular}{@{}ll@{}}
    \toprule
    Classical Stokes & Volumetric sub-problem \\
    \midrule
    velocity $u \in H^1(B; \mathbb{R}^3)$ & $~\varphi \in H^1(B; \mathbb{R}^3)$ \\
    pressure $p \in L^2(B)$ & $\pbar \in L^2(B)$ \\
    constraint $\Div u = 0$ &
      $\log J - \thbar = 0$ (via $\pbar$), $\kappa \thbar - \pbar = 0$ (via $\thbar$) \\
    pairing $\int p\,\Div v$ &
      $\int \pbar \, \F^{-\mathrm{T}}\!:\!\Grad ~\psi = \int \pbar\, \partial \theta / \partial ~\varphi [~\psi]$ \\
    \bottomrule
  \end{tabular}
  \caption{Correspondence between the classical Stokes mixed
    problem and the volumetric sub-problem of the SEC5L element.}
  \label{tab:stokes-correspondence}
\end{table}

At $\F = ~I$ the pairing reduces to the classical Stokes pairing
via $\F^{-\mathrm{T}}\!:\!\Grad ~\psi \to \Div ~\psi$, so the two
problems coincide in the small-strain limit. Away from the
reference configuration, the pairing is the nonlinear
finite-deformation generalization, but the saddle-point structure,
the role of $\pbar$ as scalar multiplier, and the inf-sup stability
requirement are the same. Consequently the stable discrete pairs
known for Stokes lift directly to the volumetric sub-problem: in
particular, the $P_k$--$P_{k-1}$ discontinuous pair is inf-sup
stable (uniformly in $h$, with a $\kappa/\mu$ dependence bounded on
the target range) and justifies the discontinuous $L^2$ choice for
$\thbar$ and $\pbar$ at degree $k-1$.

No comparable classical result exists for the deviatoric sub-problem
$(~\varphi, \ebar, \sbar)$: its stability relies on the AFW / Hu--Zhang
finite-element exterior calculus for the symmetric elasticity complex,
which is what forces the stronger $H(\Curl; \Symsd)$ and
$H(\Div; \Symsd)$ conformities there.

\textit{Determination of the polynomial degrees.} The degrees are
not free either: on trimmed (Whitney / first-kind N\'ed\'elec)
spaces each operator of the
complex~\eqref{eq:elasticity-complex} drops one polynomial degree,
so selecting continuous barycentric polynomials of degree $k$ for $~\varphi$
forces the four auxiliary slots to degree $k-1$:
\begin{equation}
  \label{eq:degree-ladder}
  \underbrace{k}_{~\varphi}
  \;\xrightarrow{\mathrm{sym}\,\Grad}\;
  \underbrace{k-1}_{\ebar}
  \;\xrightarrow{\mathrm{inc}}\;
  \underbrace{k-1}_{\sbar}
  \;\xrightarrow{\Div}\;
  \underbrace{k-1}_{\pbar,\thbar}.
\end{equation}
The scalar slots $\thbar, \pbar$ acquire degree $k-1$ because they
are dual to, respectively, the trace of $\Grad ~\varphi$ (which is
a degree-$(k-1)$ scalar) and itself. The choice $k=2$ is the
lowest robust one: $k=1$ triggers the $P_1$-$P_0$ analog of the
Stokes counterexample on tetrahedra and is generically
inf-sup-unstable for mixed elasticity without enrichment, whereas
$k=2$ yields linear-discontinuous scalar fields for $\thbar$ and
$\pbar$ with inf-sup constants bounded independently of $h$ and of
$\kappa/\mu$ on the range above.

\subsection{Discrete Residuals}
\label{sec:discrete-residuals}

Substituting~\eqref{eq:basis-functions} into~\eqref{eq:variations}
yields five discrete residual equations:
\begin{equation}
  \label{eq:discrete-residuals}
  \begin{split}
    ~R_{~\varphi;a} & :=
    \int_B \pbar \, \frac{\partial \theta}{\partial ~\varphi_a} \diff V
    +
    \int_B \sbar : \frac{\partial \edev}{\partial ~\varphi_a} \diff V
    \\
    & \quad
    -
    \int_B R ~B \, N_a \diff V
    -
    \int_{\partial_T B} ~T \, N_a \diff S
    = \mathbf{0},
    \\
    R_{\thbar;m} & :=
    \int_B \mu_m \left( \kappa \thbar - \pbar \right) \diff V
    = 0,
    \\
    ~R_{\ebar;\alpha} & :=
    \int_B ~\lambda^e_\alpha : \left( 2\mu \ebar - \sbar \right) \diff V
    = \mathbf{0},
    \\
    R_{\pbar;n} & :=
    \int_B \mu_n \left( \theta(~\varphi) - \thbar \right) \diff V
    = 0,
    \\
    ~R_{\sbar;\beta} & :=
    \int_B ~\lambda^s_\beta : \left( \edev(~\varphi) - \ebar \right) \diff V
    = \mathbf{0}.
  \end{split}
\end{equation}
The kinematic derivatives in~\eqref{eq:discrete-residuals}$_1$
follow from~\eqref{eq:kinematic-derivatives} by taking
$~\psi = N_a ~\psi_a$ and isolating the DOF direction
$~\psi_a$:
\begin{equation}
  \label{eq:kinematic-derivatives-discrete}
  \frac{\partial \theta}{\partial ~\varphi_a}
  =
  \F^{-\mathrm{T}} \Grad N_a,
  \qquad
  \frac{\partial \edev}{\partial ~\varphi_a}
  =
  \mathbb{L}\supdev(\C) : \sym\!\left(\F \otimes^s \Grad N_a\right),
\end{equation}
where $\F^{-\mathrm{T}} \Grad N_a \in \mathbb{R}^3$ is the
volumetric B-vector and
$(\F \otimes^s \Grad N_a)_{ijk} :=
\tfrac{1}{2}\left(F_{ij} N_{a,k} + F_{ik} N_{a,j}\right) \in \Syms \otimes
\mathbb{R}^3$ is the deviatoric B-tensor, both local to node $a$
and evaluated quadrature-point-wise.

\subsection{Projection Structure}
\label{sec:projection}

The four auxiliary equations~\eqref{eq:discrete-residuals}$_{2-5}$
are each $L^2$ projections from pointwise quantities onto the
corresponding finite element subspace. Define the constant mass
operators in the reference configuration:
\begin{equation}
  \label{eq:mass-matrices}
  \begin{split}
    \left[M_\theta\right]_{mn} & := \int_B \mu_m \mu_n \diff V,
    \qquad
    \left[M_e\right]_{\alpha\beta} := \int_B ~\lambda^e_\alpha : ~\lambda^e_\beta \diff V,
    \\
    \left[M_p\right]_{mn} & := \int_B \mu_m \mu_n \diff V,
    \qquad
    \left[M_s\right]_{\alpha\beta} := \int_B ~\lambda^s_\alpha : ~\lambda^s_\beta \diff V.
  \end{split}
\end{equation}
The discrete auxiliary fields are recovered from $~\varphi$ by
solving four linear projections with pre-factored constant SPD
matrices:
\begin{equation}
  \label{eq:projections}
  \begin{split}
    M_\theta \cdot \left\{\thbar\right\}_{\mathrm{dof}}
    & =
    \int_B \mu_m \log J(~\varphi) \diff V,
    \\
    M_e \cdot \left\{\ebar\right\}_{\mathrm{dof}}
    & =
    \int_B ~\lambda^e_\alpha : \edev(~\varphi) \diff V,
    \\
    M_p \cdot \left\{\pbar\right\}_{\mathrm{dof}}
    & =
    \int_B \mu_m \, \kappa \log J(~\varphi) \diff V,
    \\
    M_s \cdot \left\{\sbar\right\}_{\mathrm{dof}}
    & =
    \int_B ~\lambda^s_\alpha : \left[2\mu \, \edev(~\varphi) - 2\mu \, \edev^{\mathrm{p}}\right] \diff V.
  \end{split}
\end{equation}
In a total-Lagrangian setting the mass matrices in
\eqref{eq:mass-matrices} are independent of $~\varphi$ and of time;
a single factorization at setup (block-diagonal per element for
discontinuous $L^2$ fields, sparse and globally coupled for
conforming $H(\Curl)$ and $H(\Div)$ fields) serves every Newton
iteration of every time step.

Substituting the projected auxiliary fields
into~\eqref{eq:discrete-residuals}$_1$ collapses the five-field
system to a single discrete statement in $~\varphi$ alone:
\begin{equation}
  \label{eq:reduced-statement}
  \begin{split}
    ~R_a(~\varphi) & :=
    \int_B \pbar(~\varphi) \, \frac{\partial \theta}{\partial ~\varphi_a} \diff V
    +
    \int_B \sbar(~\varphi) : \frac{\partial \edev}{\partial ~\varphi_a} \diff V
    \\
    & \quad
    -
    \int_B R ~B \, N_a \diff V
    -
    \int_{\partial_T B} ~T \, N_a \diff S
    = \mathbf{0},
  \end{split}
\end{equation}
in which $\pbar(~\varphi)$ and $\sbar(~\varphi)$ are understood as
the evaluations of the projections~\eqref{eq:projections} applied at
the current $~\varphi$. The coupling to $\thbar(~\varphi)$ and
$\ebar(~\varphi)$ enters implicitly through the plastic
update~\eqref{eq:constitutive-plastic} and through the kinematic
dependence of $\eps(~\varphi)$ on $~\varphi$; in the reduced
statement~\eqref{eq:reduced-statement} these fields are used only
to supply the stress fields needed for equilibrium.


\section{Newton Solution}
\label{sec:newton}

\subsection{Linearization}
\label{sec:linearization}

The reduced residual~\eqref{eq:reduced-statement} is a nonlinear
function of $~\varphi$ alone. A Newton iteration computes
$\triangle ~\varphi$ from
\begin{equation}
  \label{eq:newton}
  ~R_a(~\varphi_n) + ~K_{ab} \cdot \triangle ~\varphi_b = \mathbf{0},
  \qquad
  ~\varphi_{n+1} = ~\varphi_n + \triangle ~\varphi,
\end{equation}
with tangent
\begin{equation}
  \label{eq:tangent}
  ~K_{ab}
  :=
  \frac{\partial ~R_a}{\partial ~\varphi_b}
  =
  ~K^{\mathrm{mat}}_{ab}
  +
  ~K^{\mathrm{geom}}_{ab},
\end{equation}
where the material tangent assembles contributions from both
independent stress fields through the chain rule of the projections:
\begin{equation}
  \label{eq:material-tangent}
  \begin{split}
    ~K^{\mathrm{mat}}_{ab}
    & =
    \int_B
    \frac{\partial \eps(~X)}{\partial ~\varphi_a}
    :
    \left[
      \kappa\, \mathcal{P}_{p} \, ~I \otimes ~I
      +
      2\mu \, \mathcal{P}_{s} \, \mathcal{P}_{\dev}
    \right]
    :
    \frac{\partial \eps(~X)}{\partial ~\varphi_b}
    \diff V,
  \end{split}
\end{equation}
in which $\mathcal{P}_p$ and $\mathcal{P}_s$ denote the $L^2$
projection operators onto $Q_{p,h}$ and $V_{s,h}$ respectively, and
$\mathcal{P}_{\dev}$ is the algebraic projector onto deviatoric
tensors. The geometric tangent $~K^{\mathrm{geom}}_{ab}$ arises from
the nonlinear dependence of $\eps$ on $\F$ through the matrix
logarithm and is obtained analytically via the Fr\'echet derivative
$\mathbb{L}(\C) = \partial \log\C / \partial \C$, evaluated by the
hybrid algorithm of \citet{Ortiz.Radovitzky.Repetto:2001}.

The effective tangent is pre-factorized at setup (for the parts of
the projections that are constant in the reference configuration)
and assembled at each Newton iteration. The global nonlinear system
solved in~\eqref{eq:newton} is sized by the displacement degrees of
freedom only.

\subsection{Dynamic Extension}
\label{sec:dynamics}

For transient analyses, the momentum balance in
\eqref{eq:euler-lagrange}$_1$ is augmented by the inertial term:
\begin{equation}
  \label{eq:dynamic-momentum}
  R \, \ddot{~\varphi} - \Div \P = R ~B.
\end{equation}
The auxiliary fields $\thbar$, $\ebar$, $\pbar$, $\sbar$ carry no
time derivatives: their weak statements are algebraic constraints at
each instant. Consequently the mass matrix
\begin{equation}
  \label{eq:mass}
  \left[M_{~\varphi}\right]_{ab}
  :=
  \int_B R \, N_a N_b \diff V
\end{equation}
is concentrated in the displacement block, and the semi-discrete
system is a differential-algebraic system of index one. Standard
Newmark-$\beta$ or Hilber-Hughes-Taylor integration applied to the
reduced equation~\eqref{eq:reduced-statement} with the mass
term~\eqref{eq:mass} proceeds identically as in single-field
dynamics. The auxiliary-field projections~\eqref{eq:projections} are
executed at each quadrature point evaluation of the residual and
tangent, using the current $~\varphi$.


\section{Plasticity Extension}
\label{sec:plasticity}

For isotropic $J_2$ plasticity with isotropic hardening, the
Helmholtz free energy~\eqref{eq:W-split} is modified to depend on
the elastic deviatoric log-strain
$\edev^{\mathrm{e}} := \ebar - \edev^{\mathrm{p}}$
and the equivalent plastic strain $\bar\epsilon^{\mathrm{p}}$:
\begin{equation}
  \label{eq:W-plastic}
  A(\thbar, \ebar, \edev^{\mathrm{p}}, \bar\epsilon^{\mathrm{p}})
  =
  \tfrac{\kappa}{2} \thbar^2
  +
  \mu \norm{\ebar - \edev^{\mathrm{p}}}^2
  +
  \mathcal{H}(\bar\epsilon^{\mathrm{p}}),
\end{equation}
with $\mathcal{H}$ the hardening potential. The yield function is
\begin{equation}
  \label{eq:yield}
  f(\sbar, \bar\epsilon^{\mathrm{p}})
  :=
  \norm{\sbar}
  -
  \sqrt{\tfrac{2}{3}} \left[\sigma_y + \mathcal{H}'(\bar\epsilon^{\mathrm{p}})\right].
\end{equation}
The return map is applied at each quadrature point between the
kinematic evaluation of $\edev(~\varphi)$ and the projection onto
$V_{s,h}$, reproducing the small-strain radial return of
\citet{Simo:1992}:
\begin{equation}
  \label{eq:return-map}
  \begin{split}
    \edev^{\mathrm{trial}}
    & := \edev(~\varphi)
        - \edev^{\mathrm{p}}_{n},
    \\
    \sdev^{\mathrm{trial}}
    & := 2\mu \, \edev^{\mathrm{trial}},
    \\
    \triangle \gamma
    & :=
    \frac{\norm{\sdev^{\mathrm{trial}}} - \sqrt{2/3}\,[\sigma_y + \mathcal{H}'(\bar\epsilon^{\mathrm{p}}_{n})]}
         {2\mu + \tfrac{2}{3} \mathcal{H}''(\bar\epsilon^{\mathrm{p}}_{n})}
    \quad\text{if } f^{\mathrm{trial}} > 0,
    \quad \triangle\gamma := 0 \text{ otherwise},
    \\
    ~n & := \sdev^{\mathrm{trial}} / \norm{\sdev^{\mathrm{trial}}},
    \\
    \edev^{\mathrm{p}}_{n+1}
    & := \edev^{\mathrm{p}}_{n} + \triangle \gamma \, ~n,
    \\
    \bar\epsilon^{\mathrm{p}}_{n+1}
    & := \bar\epsilon^{\mathrm{p}}_{n} + \sqrt{\tfrac{2}{3}}\,\triangle \gamma,
    \\
    \sdev^{\mathrm{pointwise}}
    & := 2\mu \left( \edev(~\varphi) - \edev^{\mathrm{p}}_{n+1} \right).
  \end{split}
\end{equation}
The pointwise updated $\sdev^{\mathrm{pointwise}}$ is then
projected onto $V_{s,h}$ via~\eqref{eq:projections}$_4$. Plastic
internal variables $\edev^{\mathrm{p}}$ and
$\bar\epsilon^{\mathrm{p}}$ are stored at quadrature points and
evolved pointwise; they are never projected, preserving the
exactness of the return map. The Lie-group/log-exp interpolation
scheme of \citet{Mota.etal:2013} applies when the internal variables
themselves are tensorial quantities requiring frame-indifferent
interpolation across quadrature.


\section{Discussion}
\label{sec:discussion}

The SEC5L element inherits three structural properties from its
symmetric-elasticity-complex foundation that distinguish it from the
composite tetrahedron:
\begin{enumerate}
  \item \textbf{Pointwise discrete symmetry of stress.} The
    independent stress field $\sbar$ lives in
    $H(\Div; \Symsd)$, so its symmetry is a property of the finite
    element space rather than of the constitutive law. This
    eliminates the cross-space aliasing between strain and stress
    fields that was observed to corrupt the equivalent plastic
    strain in stabilized displacement formulations.
  \item \textbf{Normal-traction continuity.} The $H(\Div)$
    conformity of $\sbar$ ensures that the normal component of
    the traction field is continuous across element faces. In
    regimes with large bulk-to-shear ratio, this continuity
    suppresses pressure checkerboarding that would otherwise
    develop with discontinuous stress representations.
  \item \textbf{Algebraically linear constitutive relation.} The
    Hencky--Kirchhoff pairing renders the constitutive law
    \eqref{eq:constitutive} linear in log-strain space. The
    $J_2$ return map~\eqref{eq:return-map} is identical in form to
    the small-strain case of \citet{Simo:1992}; no iteration is
    required beyond the closed-form projection. All
    finite-deformation nonlinearity is confined to the kinematic
    evaluation $\eps = \tfrac{1}{2}\log(\transpose{\F}\F)$ and the
    push-forward in~\eqref{eq:P-reconstruction}.
\end{enumerate}
The projection-and-eliminate architecture combines these structural
properties with a linear-algebra cost commensurate with a
single-field displacement formulation. Pre-factorization of the
four constant mass matrices~\eqref{eq:mass-matrices} at setup
amortizes their cost over all Newton iterations of all time steps;
each iteration incurs only four forward-solves (fast against
pre-factored SPD operators) beyond the dominant displacement Newton
solve.

Future work will address the implementation of
\begin{itemize}
  \item symmetric-tensor-valued conforming bases on tetrahedra in the
    \texttt{ReferenceFiniteElements.jl} library
    \citep{Cthonios:RFE};
  \item multi-field assembly with discontinuous $L^2$ and conforming
    $H(\Curl)$, $H(\Div)$ spaces in
    \texttt{FiniteElementContainers.jl} \citep{Cthonios:FEC};
  \item verification on canonical problems on which the composite
    tetrahedron is known to exhibit degraded performance, including
    a three-dimensional Cook's membrane at large deformation,
    contact-driven punch indentation, large-rotation twist of a
    prismatic bar, and the spherical cavity expansion benchmark
    appropriate to the high-bulk-modulus regime.
\end{itemize}
Verification of the inf-sup stability constants on these meshes is
a prerequisite for quantitative assessment. Extension to
incrementally objective integration of anisotropic plastic response
through the multiplicative decomposition $\F = \F^{\mathrm{e}} \F^{\mathrm{p}}$,
in which the Hencky--Kirchhoff conjugacy is no longer exact, is
deferred to a subsequent formulation on the gradient elasticity
complex with independent non-symmetric displacement-gradient and
first-Piola-Kirchhoff-stress fields.


\appendix

\section{Variations of the Kinematic Relation}
\label{app:variations}

This appendix makes explicit the directional derivatives
$\partial \theta / \partial ~\varphi$ and
$\partial \edev / \partial ~\varphi$ appearing in the displacement
variation~\eqref{eq:variations}$_1$. The derivation is a chain of
elementary steps expressed through the operators $\trace$, $\vol$,
and $\dev$, with only one non-elementary ingredient: the
fourth-order Fr\'echet derivative of the matrix logarithm
$\mathbb{L}(\C) := \partial \log \C / \partial \C$, whose evaluation
by the hybrid Taylor--spectral algorithm of
\citet{Ortiz.Radovitzky.Repetto:2001} is the subject of a separate
kernel that is here treated as a primitive.

\subsection{Chain of Partial Derivatives}
\label{app:chain}

The kinematic relation
$\eps(~\varphi) = \tfrac{1}{2} \log(\transpose{\F} \F)$ composes four
maps:
\begin{equation}
  \label{eq:app-composition}
  ~\varphi
  \;\stackrel{\Grad}{\longmapsto}\;
  \F
  \;\stackrel{\transpose{(\cdot)}(\cdot)}{\longmapsto}\;
  \C
  \;\stackrel{\tfrac{1}{2}\log}{\longmapsto}\;
  \eps
  \;\stackrel{\{\trace,\,\dev\}}{\longmapsto}\;
  (\theta, \edev).
\end{equation}
Applied to a test function $~\psi \in U$ with
$~\psi|_{\partial_\varphi B} = \mathbf{0}$, each link contributes a
directional derivative:
\begin{subequations}
  \label{eq:app-links}
  \begin{align}
    D_{~\varphi}\F[~\psi]
      & = \Grad ~\psi,
      \label{eq:app-dF}
    \\
    D_{\F}\C[\Grad ~\psi]
      & = \transpose{(\Grad ~\psi)}\F + \transpose{\F}\Grad ~\psi
        = 2 \sym\!\left(\transpose{\F}\Grad ~\psi\right),
      \label{eq:app-dC}
    \\
    D_{\C}\eps[\delta \C]
      & = \tfrac{1}{2}\,\mathbb{L}(\C) : \delta \C.
      \label{eq:app-deps}
  \end{align}
\end{subequations}
Both minor symmetries of $\mathbb{L}(\C)$ follow from symmetry of
$\C$; consequently
$\mathbb{L}(\C) : A = \mathbb{L}(\C) : \sym A$ for any $A$.

Composing~\eqref{eq:app-dF}--\eqref{eq:app-deps} gives the total
directional derivative of the Hencky strain:
\begin{equation}
  \label{eq:app-deps-total}
  \frac{\partial \eps}{\partial ~\varphi}[~\psi]
  =
  \mathbb{L}(\C) : \sym\!\left(\transpose{\F}\Grad ~\psi\right).
\end{equation}
Its volumetric and deviatoric parts are then obtained by the trace
and deviator projections.

\subsection{Trace Identity for the Log-Derivative}
\label{app:trace-identity}

The volumetric projection of the strain derivative reduces to a
classical identity that does not require any eigenvalue information
about $\C$. We derive it here in three steps so that the appearance
of $\C^{-1}$ (and, through $\F\C^{-1} = \F^{-\mathrm{T}}$, of
$\F^{-\mathrm{T}}$) in the volumetric variation is transparent.

\textit{Jacobi's formula.} For any invertible second-order tensor
$~A$, the derivative of the determinant with respect to its
argument is
\begin{equation}
  \label{eq:app-jacobi}
  \frac{\partial \det ~A}{\partial ~A}
  =
  (\det ~A)\,~A^{-\mathrm{T}}.
\end{equation}
Applying the chain rule to the composition $~A \mapsto \det ~A
\mapsto \log \det ~A$, the outer derivative
$\mathrm{d}(\log x)/\mathrm{d}x = 1/x$ evaluated at $x = \det ~A$
contributes a factor $1/\det ~A$ that cancels the $\det ~A$
in~\eqref{eq:app-jacobi}:
\begin{equation}
  \label{eq:app-jacobi-log}
  \frac{\partial \log \det ~A}{\partial ~A}
  =
  \frac{1}{\det ~A}\,\frac{\partial \det ~A}{\partial ~A}
  =
  \frac{1}{\det ~A}\,(\det ~A)\,~A^{-\mathrm{T}}
  =
  ~A^{-\mathrm{T}}.
\end{equation}
Specializing to the symmetric positive-definite $\C$, where
$\C^{-\mathrm{T}} = \C^{-1}$, gives the form used below:
\begin{equation}
  \label{eq:app-jacobi-C}
  \frac{\partial \log \det \C}{\partial \C}
  =
  \C^{-1}.
\end{equation}

\textit{Trace--logarithm identity.} For any symmetric positive-definite
$\C$,
\begin{equation}
  \label{eq:app-trlog}
  \trace \log \C
  =
  \log \det \C.
\end{equation}
This follows from the spectral representation: if
$\C = \sum_{i} \lambda_{i}\, ~N_{i} \otimes ~N_{i}$ with
$\lambda_{i} > 0$, then
$\trace \log \C = \sum_{i} \log \lambda_{i}
 = \log \prod_{i} \lambda_{i} = \log \det \C$.
The identity is independent of eigenvalue multiplicities and thus
robust to coalescence.

\textit{Contraction of $\mathbb{L}(\C)$ with the identity.} By
definition $\mathbb{L}(\C) = \partial \log \C / \partial \C$, and
differentiation commutes with the (constant-coefficient) trace
operator, so
\begin{equation}
  \label{eq:app-trL-step}
  ~I : \mathbb{L}(\C)
  =
  \trace \left( \frac{\partial \log \C}{\partial \C} \right)
  =
  \frac{\partial\, \trace \log \C}{\partial \C}.
\end{equation}
Combining~\eqref{eq:app-trL-step} with~\eqref{eq:app-trlog}
and~\eqref{eq:app-jacobi-C} yields the trace identity
\begin{equation}
  \label{eq:app-trL}
  ~I : \mathbb{L}(\C)
  =
  \frac{\partial \trace \log \C}{\partial \C}
  =
  \frac{\partial \log \det \C}{\partial \C}
  =
  \C^{-1}.
\end{equation}
Equation~\eqref{eq:app-trL} is a structural property of the matrix
logarithm: it holds pointwise for every symmetric positive-definite
$\C$, independently of the particular algorithm used to evaluate
$\mathbb{L}(\C)$. In particular, the hybrid Taylor--spectral scheme
of~\citet{Ortiz.Radovitzky.Repetto:2001} satisfies this identity
exactly in both of its branches, so the volumetric part of the
variation is always recovered in closed form as $\C^{-1}$ and does
not inherit any of the ill-conditioning associated with eigenvalue
coalescence.

\subsection{Volumetric Part}
\label{app:volumetric}

Applying $\trace$ to~\eqref{eq:app-deps-total} and
using~\eqref{eq:app-trL},
\begin{equation}
  \label{eq:app-dtheta}
  \begin{split}
  \frac{\partial \theta}{\partial ~\varphi}[~\psi]
  & =
  \trace\!\left\{
    \frac{\partial \eps}{\partial ~\varphi}[~\psi]
  \right\}
  =
  \left[~I : \mathbb{L}(\C)\right]
  :
  \sym\!\left(\transpose{\F}\Grad ~\psi\right)
  \\
  & =
  \C^{-1} : \sym\!\left(\transpose{\F}\Grad ~\psi\right)
  =
  \C^{-1} : \transpose{\F}\Grad ~\psi
  =
  \F^{-\mathrm{T}} : \Grad ~\psi,
  \end{split}
\end{equation}
where the penultimate step drops the $\sym$ because $\C^{-1}$ is
itself symmetric. The last equality follows from moving
$\transpose{\F}$ across the double contraction and then simplifying
$\F \C^{-1}$. Using the standard identity
$~A : (~B \, ~C) = (\transpose{~B} \, ~A) : ~C$ with
$~A = \C^{-1}$, $~B = \transpose{\F}$, and
$~C = \Grad ~\psi$,
\begin{equation}
  \label{eq:app-dtheta-step}
  \C^{-1} : \transpose{\F}\Grad ~\psi
  =
  \left(\F \C^{-1}\right) : \Grad ~\psi,
\end{equation}
and the product $\F \C^{-1}$ collapses by direct algebra:
\begin{equation}
  \label{eq:app-FCinv}
  \F \C^{-1}
  =
  \F \left(\transpose{\F}\F\right)^{-1}
  =
  \F \, \F^{-1} \F^{-\mathrm{T}}
  =
  \F^{-\mathrm{T}}.
\end{equation}
Substituting~\eqref{eq:app-FCinv} into~\eqref{eq:app-dtheta-step}
gives the last equality in~\eqref{eq:app-dtheta}. The volumetric
tensor part of the strain variation is then
$\vol(\delta \eps) = \tfrac{1}{3}\,\delta \theta\, ~I$.

Equation~\eqref{eq:app-dtheta} recovers the classical identity
$\partial \log J / \partial \F = \F^{-\mathrm{T}}$ from within the
log-strain framework. In the reduced residual, this means that the
pressure term $\int_B \pbar\,(\partial \theta / \partial ~\varphi)$
reduces to $\int_B \pbar\, \F^{-\mathrm{T}} : \Grad ~\psi$, a form
that requires neither the full $\mathbb{L}(\C)$ nor its spectral
decomposition.

\subsection{Deviatoric Part}
\label{app:deviatoric}

Applying $\dev = \mathrm{id} - \vol$ to~\eqref{eq:app-deps-total},
\begin{equation}
  \label{eq:app-de}
  \begin{split}
  \frac{\partial \edev}{\partial ~\varphi}[~\psi]
  & =
  \dev\!\left\{
    \frac{\partial \eps}{\partial ~\varphi}[~\psi]
  \right\}
  =
  \frac{\partial \eps}{\partial ~\varphi}[~\psi]
  -
  \tfrac{1}{3}\,
  \frac{\partial \theta}{\partial ~\varphi}[~\psi]\, ~I
  \\
  & =
  \mathbb{L}(\C) : \sym\!\left(\transpose{\F}\Grad ~\psi\right)
  -
  \tfrac{1}{3}\left(\F^{-\mathrm{T}} : \Grad ~\psi\right)~I.
  \end{split}
\end{equation}
Equivalently, introducing the deviatoric log-derivative
\begin{equation}
  \label{eq:app-Ldev}
  \mathbb{L}\supdev(\C)
  :=
  \mathbb{L}(\C) - \tfrac{1}{3}\, ~I \otimes \C^{-1},
\end{equation}
the deviatoric variation has the one-line form
\begin{equation}
  \label{eq:app-de-compact}
  \frac{\partial \edev}{\partial ~\varphi}[~\psi]
  =
  \mathbb{L}\supdev(\C) : \sym\!\left(\transpose{\F}\Grad ~\psi\right).
\end{equation}
By construction $~I : \mathbb{L}\supdev(\C) = \mathbf{0}$, so the
range of~\eqref{eq:app-de-compact} is trace-free, matching the space
$\Symsd$ of $\edev$.

\textit{No further analytical reduction.} The contrast between the
volumetric and deviatoric branches is structural. In~\ref{app:volumetric},
the fourth-order $\mathbb{L}(\C)$ never appears in the final
expression: the volumetric projection is the contraction
$~I : \mathbb{L}(\C)$, and that contraction coincides with
$\partial \log \det \C / \partial \C$, which has the closed form
$\C^{-1}$ by Jacobi's formula~\eqref{eq:app-jacobi-log}. The
volumetric part thus bypasses the eigenstructure of $\C$ entirely.
In~\eqref{eq:app-de-compact}, by contrast, the deviatoric projection
only removes from $\mathbb{L}(\C)$ the single piece already known in
closed form; the remainder of $\mathbb{L}(\C)$ encodes the full
fourth-order information of the matrix-logarithm derivative and has
no closed form in general, since it depends on the divided
differences $(\log \lambda_i - \log \lambda_j)/(\lambda_i - \lambda_j)$
of the eigenvalues of $\C$. The deviatoric branch therefore always
requires a dedicated evaluation of $\mathbb{L}(\C)$ ---~e.g.\ the
hybrid Taylor--spectral algorithm
of~\citet{Ortiz.Radovitzky.Repetto:2001} ---~and is the branch that
inherits the eigenvalue-coalescence issue. This asymmetry motivates
handling the volumetric and deviatoric parts with different
quadrature-point routines rather than a single monolithic evaluation
of $\mathbb{L}(\C)$.

An eigenvalue-free alternative is the Pad\'e--Fr\'echet algorithm
of~\citet{AlMohy.Higham.Relton:2013}, which extends the inverse
scaling-and-squaring Pad\'e framework already invoked
in~\eqref{eq:hencky} for $\log\C$ to produce the Fr\'echet
derivative $L_{\log}(\C, ~E)$ in the same sweep, propagating the
direction $~E$ through the Sylvester equations of the square-root
chain and through the Pad\'e kernel. Applied to each of the six
symmetric basis directions on $\Syms$, the algorithm assembles the
full fourth-order $\mathbb{L}(\C)$ without any eigendecomposition,
removing the coalescence branch entirely and unifying the numerical
treatment of $\log\C$ and $\partial \log\C / \partial \C$ under a
single kernel.

\subsection{Computational Summary}
\label{app:summary}

At each quadrature point of the residual
assembly~\eqref{eq:discrete-residuals}$_1$ the evaluation proceeds in
five steps:
\begin{enumerate}
  \item Kinematics: $\F = \Grad ~\varphi_h$,
    $\C = \transpose{\F}\F$, $\F^{-\mathrm{T}}$.
  \item (Optionally) spectral decomposition of $\C$, skipped when
    $\norm{\C - ~I}$ falls below the Taylor threshold of
    \citet{Ortiz.Radovitzky.Repetto:2001}.
  \item Evaluation of $\mathbb{L}(\C)$ by the hybrid algorithm
    (spectral branch or Taylor recurrence).
  \item For each displacement basis $N_a$ and direction
    $~e_i$, with $G_{a,i} := \Grad(N_a\, ~e_i)$:
    \begin{align*}
      \delta\theta_{a,i}
        & =
        \F^{-\mathrm{T}} : G_{a,i},
      \\
      \delta\eps_{a,i}
        & =
        \mathbb{L}(\C) : \sym\!\left(\transpose{\F}\, G_{a,i}\right),
      \\
      \delta\edev_{a,i}
        & =
        \delta\eps_{a,i}
        -
        \tfrac{1}{3}\,\delta\theta_{a,i}\, ~I.
    \end{align*}
  \item Contraction with the stress fields:
    $(\pbar\, ~I + \sbar) : \delta\eps_{a,i}
    = \pbar\,\delta\theta_{a,i} + \sbar : \delta\edev_{a,i}$,
    consistent with the orthogonal
    split~\eqref{eq:voldev-strain}--\eqref{eq:voldev-stress}.
\end{enumerate}
The volumetric and deviatoric variations share a single evaluation
of $\mathbb{L}(\C)$; no separate constitutive-level call is required
for the trace, which is obtained from $\F^{-\mathrm{T}}$ alone. The
only object that depends on the kinematic state at the current
Newton iterate is $\mathbb{L}(\C)$ itself.


\section{The Symmetric Elasticity Complex}
\label{app:sec}

This appendix collects the structural background behind the
complex~\eqref{eq:elasticity-complex}. The goal is operational: to
identify each operator, its kernel, its range, and its interaction
with the orthogonal volumetric--deviatoric split
$\Syms = \Syms\supvol \oplus \Symsd$ that organizes the five-field
formulation.

\subsection{The de Rham Complex and its Tensorial Lift}
\label{app:sec-derham}

On a bounded, simply-connected, Lipschitz domain
$B \subset \Reals^3$, the scalar de Rham complex
\begin{equation}
  \label{eq:app-derham}
  \Reals
  \;\xrightarrow{\iota}\;
  H^1(B)
  \;\xrightarrow{\Grad}\;
  H(\Curl; B)
  \;\xrightarrow{\Curl}\;
  H(\Div; B)
  \;\xrightarrow{\Div}\;
  L^2(B)
  \;\xrightarrow{0}\; 0
\end{equation}
is exact: $\Curl \circ \Grad = 0$, $\Div \circ \Curl = 0$, and the
kernel of each operator coincides with the range of its predecessor.
Exactness encodes the classical vector-calculus identities
(``$\Curl$-free implies gradient,'' ``$\Div$-free implies curl'') as
statements about the kernels and ranges of differential operators
between Sobolev spaces.

The symmetric elasticity complex is the symmetric-tensor analog
of~\eqref{eq:app-derham}. The Bernstein--Gel'fand--Gel'fand (BGG)
construction of~\citet{Arnold.Falk.Winther:2006a,
Arnold.Falk.Winther:2010} couples a scalar de Rham complex with a
vector-valued one through an algebraic projection, resolving the
resulting double complex to yield
\begin{equation}
  \label{eq:app-sec}
  \operatorname{RM}(B)
  \;\xrightarrow{\iota}\;
  H^1(B; \Reals^3)
  \;\xrightarrow{\sym\Grad}\;
  H(\Curl; B; \Syms)
  \;\xrightarrow{\mathrm{inc}}\;
  H(\Div; B; \Syms)
  \;\xrightarrow{\Div}\;
  L^2(B; \Reals^3)
  \;\xrightarrow{0}\; 0,
\end{equation}
where $\operatorname{RM}(B) = \{ ~a + ~W ~X : ~a \in \Reals^3,\,
~W \in \mathrm{skew}(\Reals^3) \}$ is the six-dimensional space of
infinitesimal rigid motions. Each space in~\eqref{eq:app-sec} carries
a symmetric-tensor (or vector) field with Sobolev conformity
dictated by the operator it must survive; the complex is exact on a
simply-connected domain.

\subsection{Operators, Kernels, and Ranges}
\label{app:sec-ops}

The three differential operators of~\eqref{eq:app-sec} are
\begin{enumerate}
  \item \textbf{Symmetric gradient.}
    $\sym\Grad : H^1(B; \Reals^3) \to L^2(B; \Syms)$ with
    \begin{equation}
      \label{eq:app-symgrad}
      \sym\Grad ~u
      :=
      \tfrac{1}{2}\!\left(\Grad ~u + \transpose{\Grad ~u}\right).
    \end{equation}
    Its kernel is $\operatorname{RM}(B)$, the space of infinitesimal
    rigid motions; this is the Korn--Kirchhoff theorem.
  \item \textbf{Incompatibility.}
    $\mathrm{inc} : H(\Curl; B; \Syms) \to H(\Div; B; \Syms)$ is the
    second-order operator
    \begin{equation}
      \label{eq:app-inc}
      \mathrm{inc}(~\epsilon)
      :=
      \Curl\,\transpose{(\Curl\,~\epsilon)},
      \qquad
      \left[\mathrm{inc}(~\epsilon)\right]_{pq}
      =
      \varepsilon_{pij}\,\varepsilon_{qkl}\,\epsilon_{jl,ik},
    \end{equation}
    where the $\Curl$ of a tensor field acts row-wise and
    $\varepsilon_{ijk}$ is the Levi-Civita symbol. The kernel of
    $\mathrm{inc}$ coincides with the range of $\sym\Grad$: a
    symmetric-tensor field $~\epsilon$ is compatible with a
    displacement, $~\epsilon = \sym\Grad ~u$, if and only if
    $\mathrm{inc}(~\epsilon) = 0$. This is the classical
    \emph{Saint-Venant compatibility condition}.
  \item \textbf{Divergence.}
    $\Div : H(\Div; B; \Syms) \to L^2(B; \Reals^3)$ acts row-wise,
    $(\Div ~\tau)_i = \tau_{ij,j}$. Its kernel is the range of
    $\mathrm{inc}$: a symmetric-tensor stress field is self-equilibrated
    if and only if it is the $\mathrm{inc}$ of a symmetric-tensor
    strain field.
\end{enumerate}
The complex loses one order of regularity at each step
($H^1 \to H(\Curl) \to H(\Div) \to L^2$) and is exact on a
simply-connected domain. The exactness at the $L^2$ slot is the
solvability of the equilibrium equation $\Div~\tau = -R~B$ in
$H(\Div; \Syms)$ for any admissible body-force data.

\subsection{Interaction with the Volumetric--Deviatoric Split}
\label{app:sec-voldev}

The orthogonal decomposition $\Syms = \Syms\supvol \oplus \Symsd$
does \emph{not} descend to a pair of sub-complexes of~\eqref{eq:app-sec}.
The middle operator $\mathrm{inc}$ mixes volumetric and deviatoric
parts, as the following explicit computation shows.

Take a purely volumetric input $~\epsilon = \alpha ~I$ with
$\alpha \in H^1(B)$ and $~I$ the second-order identity. We regard
$\alpha~I$ as an element of $H(\Curl; \Syms)$, not as a compatible
strain: for generic non-affine $\alpha$ the field $\alpha~I$ lies
outside the range of $\sym\Grad$, since
$\mathrm{inc}(\alpha~I) = 0$ forces $\alpha_{,pq} = 0$ for $p \neq q$
and $\Delta\alpha = \alpha_{,pp}$ (no sum), whose only solutions are
affine $\alpha = a + b_i X_i$. The counterexample below is
informative precisely because it probes $\mathrm{inc}$ on the
non-compatible part of the middle space. Substituting
$\epsilon_{jl} = \alpha\,\delta_{jl}$ into~\eqref{eq:app-inc} and
contracting with the Levi-Civita identity
$\varepsilon_{pij}\,\varepsilon_{qkj} = \delta_{pq}\delta_{ik} -
\delta_{pk}\delta_{iq}$,
\begin{equation}
  \label{eq:app-inc-vol}
  \left[\mathrm{inc}(\alpha~I)\right]_{pq}
  =
  \varepsilon_{pij}\,\varepsilon_{qki}\,\alpha_{,jk}
  =
  \delta_{pq}\,\alpha_{,ii} - \alpha_{,pq},
\end{equation}
or in coordinate-free form,
\begin{equation}
  \label{eq:app-inc-vol-cf}
  \mathrm{inc}(\alpha ~I)
  = (\Delta \alpha)\, ~I - \Grad\Grad\alpha.
\end{equation}
Projecting onto the vol--dev split,
\begin{equation}
  \label{eq:app-inc-vol-split}
  \mathrm{inc}(\alpha ~I)
  = \underbrace{\tfrac{2}{3}(\Delta\alpha)\,~I}_{\vol}
  \;-\;
  \underbrace{\!\left[\Grad\Grad\alpha - \tfrac{1}{3}(\Delta\alpha)\,~I\right]\!}_{\dev},
\end{equation}
where the deviatoric part is the trace-free Hessian of $\alpha$,
generically nonzero. A purely volumetric strain is therefore mapped
by $\mathrm{inc}$ to a stress with nonzero deviatoric part. The
converse also fails: the trace of $\mathrm{inc}\,~\epsilon$ is
$\trace[\mathrm{inc}\,~\epsilon] = \epsilon_{ij,ij}$, which is
nonzero for a generic deviatoric $~\epsilon \in \Symsd$.

This is the structural reason why the volumetric and deviatoric
sub-problems of the five-field formulation cannot be decoupled at the
level of the complex itself: the complex-compatible conformities
$H(\Curl; \Syms)$ and $H(\Div; \Syms)$ act on full symmetric tensors,
and the trace-free restrictions inherit their stability from the
full-tensor theory rather than from an independent sub-complex
result. What the vol--dev split \emph{does} provide is an
orthogonal decomposition of the constitutive operator through its
action on a strain field,
\begin{equation}
  \label{eq:app-Cvoldev}
  \mathbb{C}\!:\!~\epsilon
  =
  \kappa\,(\trace ~\epsilon)\,~I
  +
  2\mu\,\dev ~\epsilon,
\end{equation}
which decouples the bulk and shear responses into additive terms
whose ranges lie in $\Syms\supvol$ and $\Symsd$ respectively. This
decoupling is what allows the scalar fields $\thbar$ and $\pbar$ to
be extracted as trace-projections of the tensor fields and promoted
to independent variables with their own inf-sup condition. The deviatoric tensor fields $\ebar \in H(\Curl; \Symsd)$
and $\sbar \in H(\Div; \Symsd)$ inherit conformity from the
trace-free projection of the full SEC spaces.

\subsection{Discrete Spaces and Commuting Projections}
\label{app:sec-discrete}

The finite-dimensional analog of~\eqref{eq:app-sec} is a sequence of
conforming spaces $V_{~\varphi,h}$, $V_{\ebar,h}$, $V_{\sbar,h}$,
$V_{\pbar,h}$ on a tetrahedral mesh together with interpolation
operators $\Pi_{~\varphi}, \Pi_{\ebar}, \Pi_{\sbar}, \Pi_{\pbar}$
that commute with the differential operators. For sufficiently
smooth fields,
\begin{equation}
  \label{eq:app-commuting}
  \begin{CD}
    H^1(B; \Reals^3)
    @>\sym\Grad>>
    H(\Curl; \Syms)
    @>\mathrm{inc}>>
    H(\Div; \Syms)
    @>\Div>>
    L^2(B; \Reals^3)
    \\
    @V \Pi_{~\varphi} VV
    @V \Pi_{\ebar} VV
    @V \Pi_{\sbar} VV
    @V \Pi_{\pbar} VV
    \\
    V_{~\varphi,h}
    @>\sym\Grad>>
    V_{\ebar,h}
    @>\mathrm{inc}>>
    V_{\sbar,h}
    @>\Div>>
    V_{\pbar,h}.
  \end{CD}
\end{equation}
The commuting diagram delivers two structural consequences: the
discrete complex is itself exact (modulo the kernel of
$\Pi_{~\varphi}$), and the inf-sup stability of the coupled system
can be read off the continuous complex by pulling stability through
the projections. This is the central result of finite-element
exterior calculus~\citep{Arnold.Falk.Winther:2010}: cochain
projections that commute with the differential operators transfer
stability from the continuous to the discrete level.

Two families of discrete symmetric-tensor spaces realize the commuting
diagram~\eqref{eq:app-commuting} on tetrahedra.

\begin{enumerate}
  \item \textbf{Strongly-symmetric elements.} The
    Arnold--Awanou--Winther
    family~\citep{Arnold.Awanou.Winther:2008} enforces symmetry of
    the tensor field pointwise through a vertex-edge-face-cell
    hierarchy of degrees of freedom. The $H(\Div; \Syms)$ element of
    lowest order on a tetrahedron uses polynomial degree $k = 3$ and
    $162$ degrees of freedom per element. The construction is
    geometrically clean --- symmetry is a property of the finite
    element space --- at the cost of a high local dimension at low
    polynomial order.
  \item \textbf{Weakly-symmetric elements.} The
    Arnold--Falk--Winther family~\citep{Arnold.Falk.Winther:2007}
    trades pointwise symmetry for a smaller local dimension by
    enforcing symmetry variationally. A skew-symmetric Lagrange
    multiplier $~W \in L^2(B; \mathbb{K})$, with $\mathbb{K}$ the
    space of skew-symmetric second-order tensors, imposes the
    constraint
    $\int_B ~W : ~\tau \diff V = 0$ on the tensor field, and
    the $H(\Div)$ space itself need not be pointwise-symmetric.
    The resulting mixed element uses substantially fewer degrees of
    freedom at low polynomial order, at the cost of adding the skew
    multiplier to the system.
\end{enumerate}

Both families admit cochain projections compatible
with~\eqref{eq:app-commuting} and are established stable
discretizations of the symmetric elasticity complex. The present
formulation inherits inf-sup stability from whichever family is
selected in \texttt{ReferenceFiniteElements.jl}
\citep{Cthonios:RFE}; the choice affects the element-local assembly
but does not alter the variational formulation~\eqref{eq:hw} or the
reduced residual~\eqref{eq:discrete-residuals}.


\bibliographystyle{plainnat}
\bibliography{references}

\end{document}
